% Integrated research report: NLP for Structured HAZOP Reasoning
\documentclass[11pt,a4paper]{report}
\usepackage[a4paper,margin=0.85in]{geometry}
\usepackage[T1]{fontenc}
\usepackage{lmodern,microtype}
\usepackage{amsmath,amssymb}
\usepackage{booktabs,tabularx,longtable,array}
\usepackage{enumitem}
\usepackage{xcolor}
\usepackage[most]{tcolorbox}
\usepackage{hyperref}
\definecolor{hzblue}{HTML}{1F4E79}
\definecolor{hzred}{HTML}{8C1D1D}
\definecolor{hzredbg}{HTML}{FBE4E4}
\hypersetup{colorlinks=true,linkcolor=hzblue,citecolor=hzblue,urlcolor=hzblue,pdftitle={Integrated Research Report: NLP for Structured HAZOP Reasoning},pdfauthor={Prasangeet Dongre}}
\newtcolorbox{boundary}[1][]{enhanced,breakable,colback=hzredbg,colframe=hzred,title={#1},fonttitle=\bfseries}
\setlist{topsep=4pt,itemsep=3pt}
\renewcommand{\arraystretch}{1.15}
\title{\textbf{NLP for Structured HAZOP Reasoning}\\[0.5em]\large Integrated Research and Experimental Report}
\author{Prasangeet Dongre\\Roll No.: B23CH1033\\Department of Chemical Engineering\\Indian Institute of Technology Jodhpur\\[0.3em]\small Advisor: Dr. Santhosh Kumar Varanasi}
\date{November 2026}
\begin{document}
\maketitle
\begin{abstract}
This report consolidates the research framing, literature review, repository architecture, dataset tracks, synthetic-data provenance, and seven experiments (E1--E7) for the NLP for Structured HAZOP Reasoning project. The system represents HAZOP information using typed relationships and evaluates recovery of held-out cause-to-consequence worksheet associations. On the fixed reconstructed T-301 split, graph reasoning achieves micro-F1 of 0.7442, compared with 0.3553 for TF--IDF and 0.5549 for the deterministic graph--TF--IDF union. Follow-on experiments show sensitivity to evidence support, scenario splits, provenance filters, and controlled evidence deletion. The target labels are reconstructed row-level silver associations, not expert-adjudicated physical causal relationships. Synthetic T-301 continuation and DWSIM scenarios are kept separate from the reconstructed benchmark.
\end{abstract}
\noindent\textbf{Keywords:} HAZOP; process safety; NLP; knowledge graphs; causal reasoning; provenance; LLMs.
\tableofcontents
\chapter{Research Problem and Scope}
\section{Background}
HAZOP worksheets record process nodes, design intent, guidewords, parameters, deviations, causes, consequences, safeguards, and recommendations. A flat worksheet row may omit intermediate events, branching pathways, and the attachment points of safeguards. The project investigates whether structured NLP and graph-based methods can assist recovery of missing relationships in HAZOP documentation.
\section{Implemented scope}
The implemented and evaluated scope covers typed HAZOP records and relationships, provenance tracking, CSV loading, graph persistence and querying, lexical and graph-neighbourhood retrieval, bounded graph reasoning, hybrid candidate generation, and benchmark evaluation on reconstructed T-301 data. The repository does not establish a fully validated pipeline from arbitrary scanned reports to expert-approved causal graphs.
\section{Research questions}
The experiments ask: (1) how retrieval and graph methods compare; (2) how minimum evidence support affects recovery; (3) whether independent edge thresholds differ from complete-path thresholds; (4) whether the dataset exercises intermediate multi-hop paths; (5) whether provenance completeness and source coherence improve proposals; (6) how results vary across scenario splits; and (7) how controlled relationship deletion affects recovery.
\begin{boundary}[Interpretation boundary]
E1--E7 predict whether a consequence is co-listed with a cause in a reconstructed worksheet row. These silver labels are not proof that the cause physically produces every listed consequence. Metrics do not establish causal correctness, expert approval, or process-safety utility.
\end{boundary}

\chapter{Literature Review}
\section{Classical model-based and expert-system approaches}
Digraph-based and expert-system approaches reason over process models or encoded expert knowledge. Representative references in the project's bibliography include Vaidhyanathan and Venkatasubramanian, ExpHAZOP+, D-higraphs, and later automated HAZOP work \cite{vaidhyanathan1995digraph,vaidhyanathan1996experience,rahman2009exphazop,rodriguez2012dhigraphs,taylor2017automated}.
\section{Ontology-based reasoning}
Ontology-based systems explicitly represent HAZOP concepts and relations. Single et al. describe an ontology-based computer aid for HAZOP studies \cite{single2020ontologyhazop}. This work motivates typed representation and symbolic reasoning, while differing from extraction of structure from arbitrary historical documents.
\section{NLP and statistical text mining}
The reviewed literature includes NLP applied to HAZOP reports, active-learning text mining, deviation prediction, and recommendation prediction \cite{feng2021nlphazop,wang2021activelearning,ekramipooya2023deviation,ekramipooya2024recommendation}. Entity or label extraction is not equivalent to recovering a directed and contextually valid causal chain.
\section{Knowledge graphs and LLM generation}
Process-safety knowledge graph construction and graph-based HAZOP completeness work motivate explicit graph representation \cite{mao2020processsafetykg,feng2021hazopcausal,wang2022iskghainex,lizhao2024constructing,lizhao2025hazopct}. Direct LLM worksheet generation is a different task from evidence-grounded relationship recovery \cite{lee2026llmhazop}.
\section{Research position}
This project investigates typed relationship representation, provenance preservation, and bounded recovery of missing worksheet associations. It does not yet provide expert-adjudicated evaluation of physical causal chains.

\chapter{Representation and Software Architecture}
\section{Typed domain model}
The domain representation includes HAZOP records and typed relationships. Relationship metadata includes subject, predicate, object, provenance origin, evidence, confidence, source identifier, endpoint types, and source location. Predicates include \texttt{causes\_deviation}, \texttt{leads\_to}, \texttt{leads\_to\_consequence}, \texttt{co\_occurs\_with\_consequence}, \texttt{co\_listed\_with\_consequence}, \texttt{mitigated\_by}, \texttt{recommends\_action}, \texttt{has\_parameter}, and \texttt{located\_in}.
\section{Provenance}
Record status distinguishes source-reconstructed and synthetic records from expert-validated records. Relationship origins distinguish source, reconstructed, synthetic, expert, model-inferred, graph-inferred, retrieval-inferred, and LLM-extracted propositions. Provenance supports traceability; it does not itself establish truth. In particular, \texttt{co\_listed\_with\_consequence} denotes a worksheet association, not a verified physical causal edge.
\section{Implementation}
The repository contains a FastAPI backend and Tauri--React client. Backend components include domain entities and ontology definitions, dataset loaders and repositories, API routes, retrieval strategies, and graph/LLM reasoning strategies. In-memory and Neo4j-oriented graph persistence are present. Benchmark results describe the specific scripts and methods evaluated; they do not prove that every planned application capability is complete.

\chapter{Datasets and Synthetic Scenarios}
\section{T-301 reconstructed and continuation tracks}
The canonical file \texttt{data/raw/t301/T301\_HAZOP\_CCPS2001\_style.csv} contains 94 records: 40 \texttt{PAPER\_REPORTED\_RECONSTRUCTION} records and 54 \texttt{SYNTHETIC\_CONTINUATION} records. The 40 records are structured reconstructions of scenario chains reported for the T-301 hexane-storage case in Single et al. (2020); they are not verbatim source worksheet rows. The 54 continuation records are project-generated and must not be attributed to the paper. E1--E7 use only the 40 reconstructed records.
\section{DWSIM synthetic track}
The canonical file \texttt{data/raw/dwsim/synthetic\_hazop\_D001\_D030.csv} contains 30 synthetic scenarios, D001--D030, associated with an ethanol/fermentation plant representation. The file records node, equipment, parameter, guideword, deviation, possible cause, simulated or inferred effect, possible consequence, candidate safeguard, evidence and limitation, source type, and validation status. The README reports 16 \emph{Review required}, 6 \emph{Scope caveat}, 6 \texttt{NEEDS\_MODEL\_REVISION}, 1 \emph{Review reaction model}, and 1 \texttt{NEEDS\_REVIEW}. The DWSIM response is captured for the simulation-related effect field; consequence and safeguard fields remain synthetic hypotheses. None of the records is thereby expert-validated.
The file \texttt{HAZOP\_DWSIM30\_worksheet\_format.csv} is an alternate worksheet-style export of the same D001--D030 scenarios, not an independent additional sample.
\section{Synthetic-generation audit}
The tracked repository documents the resulting synthetic datasets and their provenance/status fields. The repository tree inspected for this report contains the CSV outputs and benchmark/import code, but no dedicated versioned scenario-generation script or complete generation manifest specifying prompts, model versions, random seeds, sampling parameters, simulation conditions, and review procedure. Therefore the exact generation algorithm cannot be reconstructed responsibly from the tracked files. The T-301 continuation is supported as project-generated by its record type and provenance; this does not establish a reproducible generation method. The DWSIM data is explicitly synthetic and simulation/inference-derived, with its limitations recorded per row.
\begin{table}[htbp]\centering\small
\caption{Dataset tracks and their role in E1--E7}
\begin{tabularx}{\textwidth}{@{}l r X X@{}}
\toprule
Track & Records & Status & E1--E7 use\\
\midrule
T-301 reconstructed & 40 & Reconstructed from reported scenario chains; not verbatim rows & Included\\
T-301 continuation & 54 & Synthetic continuation, separate from paper reconstruction & Excluded\\
DWSIM synthetic & 30 & Synthetic and simulation/inference-derived; review caveats retained & Excluded\\
DWSIM worksheet export & 30 & Alternate export of D001--D030 & Not independent\\
\bottomrule
\end{tabularx}\end{table}
\section{Silver-label construction}
For each held-out cause, the target set contains consequences co-listed in the same source row. There is no independently adjudicated relationship-level gold set in the documented E1--E7 evaluation. These labels measure recovery of worksheet associations, not physical causation.

\chapter{Experimental Protocol}
\section{Fixed split}
The fixed benchmark uses a scenario-level 80/20 split with seed 42: 32 training scenarios and 8 held-out scenarios, producing 19 test queries. All records belonging to a scenario remain in one partition. The held-out target scenario's causal labels and graph edges are excluded from the retrieval repository. Protocol version is 1.3 and top-$k=10$. LLM methods are capped at 15 evaluated cases.
The dataset SHA-256 is \texttt{1c8713c2b6fbc9cdf2e22e64a307fdb64f073158a78679bdc23ad803a7d2725e}.
\section{Metrics}
Precision, recall, and F1 are measured against row-derived silver targets. Coverage is the fraction of eligible queries receiving at least one proposal; abstention is its complement. Selective risk is the error rate among answered queries and should be read together with coverage. Latency is implementation- and environment-dependent.

\chapter{Experimental Results E1--E7}
\section{E1: Strategy comparison}
\begin{table}[htbp]\centering\scriptsize
\caption{E1 full-test strategy comparison on reconstructed T-301}
\begin{tabularx}{\textwidth}{@{}Xrrrrrrr@{}}
\toprule
Method & N/19 & P & R & F1 & Coverage & Abstention & Latency (ms)\\
\midrule
Prompt-only LLM & 15 & 0.0000 & 0.0000 & 0.0000 & 0.0667 & 0.9333 & 16802.40\\
TF--IDF retrieval & 19 & 0.2784 & 0.4909 & 0.3553 & 0.8947 & 0.1053 & 0.51\\
Graph reasoning & 19 & 0.6486 & 0.8727 & 0.7442 & 0.8421 & 0.1579 & 1.34\\
Graph + TF--IDF union & 19 & 0.4068 & 0.8727 & 0.5549 & 0.8947 & 0.1053 & 1.85\\
Hybrid graph/vector/LLM & 15 & 0.3704 & 0.6383 & 0.4688 & 0.8667 & 0.1333 & 38415.43\\
\bottomrule
\end{tabularx}\end{table}
Graph reasoning is the strongest deterministic method by micro-F1 on the fixed split. The union increases coverage but reduces precision and F1. LLM methods did not outperform graph reasoning on the matched 15-query subset and had substantially higher latency. Full-test LLM rows have different denominators and must not be treated as fully matched comparisons.
On the matched subset, TF--IDF has P=0.3462, R=0.5745, F1=0.4320, coverage=0.8667; graph has P=0.6452, R=0.8511, F1=0.7339, coverage=0.8000; union has P=0.4396, R=0.8511, F1=0.5797, coverage=0.8667; prompt-only LLM has F1=0 and coverage=0.0667; hybrid has P=0.3704, R=0.6383, F1=0.4688, coverage=0.8667.
\section{E2: Global complete-path consensus}
Both edges of a typed cause--deviation--consequence path must be supported in the same eligible training record and process family, excluding the target scenario.
\begin{table}[htbp]\centering\small
\caption{E2 complete-path support thresholds}
\begin{tabular}{@{}rrrrrr@{}}
\toprule
Minimum support & P & R & F1 & Coverage & Abstention\\
\midrule
1 & 0.6486 & 0.8727 & 0.7442 & 0.8421 & 0.1579\\
2 & 0.6000 & 0.1091 & 0.1846 & 0.3158 & 0.6842\\
3--6 & 0 & 0 & 0 & 0 & 1\\
\bottomrule
\end{tabular}\end{table}
Requiring two supporting scenarios sharply reduces recall and coverage; thresholds of three or more yield no eligible proposals. Complete-path recurrence is sparse in this dataset.
\section{E3: Independent edge support}
\begin{table}[htbp]\centering\small
\caption{E3 independent edge-support thresholds}
\begin{tabular}{@{}lrrrrr@{}}
\toprule
Condition & P & R & F1 & Coverage & Abstention\\
\midrule
Both thresholds 1 & 0.6486 & 0.8727 & 0.7442 & 0.8421 & 0.1579\\
Cause edge $\geq 2$ & 0.5938 & 0.3455 & 0.4368 & 0.3158 & 0.6842\\
Consequence edge $\geq 2$ & 0.6094 & 0.7091 & 0.6555 & 0.8421 & 0.1579\\
Both edges $\geq 2$ & 0.6000 & 0.3273 & 0.4235 & 0.3158 & 0.6842\\
\bottomrule
\end{tabular}\end{table}
Cause-edge recurrence is the bottleneck. Raising only the consequence-edge threshold has a smaller effect. Repeated co-occurrence is not independent causal confirmation.
\section{E4: Multi-hop reasoning}
\begin{table}[htbp]\centering\small
\caption{E4 path-mode comparison}
\begin{tabular}{@{}lrrrrrr@{}}
\toprule
Path mode & Candidate paths & P & R & F1 & Coverage & Abstention\\
\midrule
Direct & 265 & 0.6486 & 0.8727 & 0.7442 & 0.8421 & 0.1579\\
Intermediate & 0 & 0 & 0 & 0 & 0 & 1\\
Both & 265 & 0.6486 & 0.8727 & 0.7442 & 0.8421 & 0.1579\\
\bottomrule
\end{tabular}\end{table}
No eligible intermediate paths were present. This is a limitation of evidence availability, not proof that multi-hop reasoning inherently fails.
\section{E5: Provenance coherence and selective acceptance}
E5-1 metadata-presence thresholds 0, 0.5, 0.75, and 1 produced identical results: P=0.6486, selective risk=0.3514, F1=0.7442, coverage=0.8421, abstention=0.1579, unsupported=0, traceable=1.0. There were 74 accepted proposals: 48 TP, 26 FP, and 7 FN. Metadata completeness did not discriminate correct from incorrect proposals.
\begin{table}[htbp]\centering\small
\caption{E5 source-coherence filters}
\begin{tabular}{@{}lrrrrrr@{}}
\toprule
Condition & P & R & F1 & Coverage & Accepted & TP/FP/FN\\
\midrule
Ungated & 0.6486 & 0.8727 & 0.7442 & 0.8421 & 74 & 48/26/7\\
Same scenario & 0.5000 & 0.1091 & 0.1791 & 0.2105 & 12 & 6/6/49\\
Same source row & 0.5000 & 0.1091 & 0.1791 & 0.2105 & 12 & 6/6/49\\
\bottomrule
\end{tabular}\end{table}
Both filters reject 62 of 74 proposals and retain 6 TP and 6 FP. They reduce coverage without improving precision. Provenance coherence supports traceability, not causal truth.
\section{E6: Scenario-split robustness}
E6 repeats scenario-level 80/20 holdouts for seeds 42--51. Each run holds out 8 of 40 scenarios; query counts range from 17 to 27. Across runs, 36 unique scenarios were held out and 4 never appeared in a test partition. These overlapping runs are repeated random holdouts, not ten-fold cross-validation.
\begin{table}[htbp]\centering\small
\caption{E6 repeated-holdout summary}
\begin{tabular}{@{}lrrrr@{}}
\toprule
Metric & Mean & SD & Min & Max\\
\midrule
Micro precision & 0.5765 & 0.0412 & 0.5294 & 0.6486\\
Micro recall & 0.8615 & 0.1034 & 0.6383 & 1.0000\\
Micro F1 & 0.6881 & 0.0479 & 0.5825 & 0.7442\\
Macro F1 & 0.6266 & 0.0973 & 0.3722 & 0.7057\\
Coverage & 0.8746 & 0.1283 & 0.5500 & 1.0000\\
Abstention & 0.1254 & 0.1283 & 0 & 0.4500\\
\bottomrule
\end{tabular}\end{table}
The results show split sensitivity. The summary is descriptive and is not an independent-sample confidence interval.
\section{E7: Controlled relationship-evidence masking}
E7 fixes the seed-42 split and randomly masks 10--40\% of 181 typed training relationships. Ten masking seeds (42--51) are used at each rate, with masks nested within each seed. Test labels are unchanged and LLM methods are disabled. The zero-mask row is the E1 baseline.
\begin{table}[htbp]\centering\small
\caption{E7 graph reasoning under random relationship masking}
\begin{tabular}{@{}rrrrr@{}}
\toprule
Masked & Precision & Recall & F1 & Coverage\\
\midrule
0\% & 0.6486 & 0.8727 & 0.7442 & 0.8421\\
10\% & 0.6454 & 0.8109 & 0.7182 & 0.8000\\
20\% & 0.6439 & 0.7382 & 0.6872 & 0.7316\\
30\% & 0.6396 & 0.6582 & 0.6469 & 0.6842\\
40\% & 0.6476 & 0.5727 & 0.6030 & 0.6158\\
\bottomrule
\end{tabular}\end{table}
As masking increases, graph recall, F1, and coverage decline while precision remains comparatively stable. The graph--TF--IDF union has P/R/F1/coverage of 0.4077/0.8145/0.5432/0.8737 at 10\%, 0.4124/0.7455/0.5305/0.8526 at 20\%, 0.4190/0.6673/0.5139/0.8263 at 30\%, and 0.4299/0.5855/0.4936/0.7737 at 40\%. TF--IDF has 0.2713/0.4182/0.3288/0.8684, 0.2770/0.3655/0.3143/0.8368, 0.2663/0.2945/0.2775/0.8053, and 0.2570/0.2291/0.2402/0.7526 at the same rates. The masking procedure is uniform random deletion of typed relationship objects, not a model of natural missingness in industrial reports.

\chapter{Discussion, Limitations, and Future Work}
\section{Supported findings}
The fixed-split results identify graph reasoning as the strongest tested deterministic method by micro-F1. Candidate union increases coverage but trades away precision and F1. Complete-path and cause-edge recurrence are sparse. Provenance-field completeness does not discriminate correctness. Strict same-scenario and same-row filters over-filter proposals. The benchmark does not exercise intermediate paths. Repeated holdouts show scenario-split sensitivity, and random relationship masking degrades graph recall, F1, and coverage.
\section{Limitations}
The primary evaluated set has 40 reconstructed scenarios, eight held-out scenarios, and 19 fixed-split queries. Labels are reconstructed row associations, not expert-adjudicated causal ground truth. E6 runs overlap. E4 has no eligible intermediate paths. E7 uses uniform random deletion. LLM evaluation is capped at 15 cases. Synthetic generation is not fully reproducible from the tracked files because a complete generation script or manifest was not found in the repository tree inspected for this report.
\section{Future work}
The most important next step is an independently adjudicated relationship dataset with a written rubric for causal plausibility, context sufficiency, temporal ordering, operating envelope, and safeguard validity. Multiple qualified reviewers should annotate independently, with disagreement resolution and inter-rater agreement reported. Further work should preserve separate reconstructed, synthetic, and expert-validated tracks; document synthetic generation with prompts, model/version, parameters, seeds, simulation conditions, and review outcomes; evaluate additional process units; create examples that exercise intermediate paths; assess calibration and abstention against expert judgement; validate the full document-extraction pipeline; and test process-safety constraints with domain experts.
\begin{boundary}[Process-safety disclaimer]
The prototype is a research aid. Recovered links are candidate worksheet associations, not verified process-safety facts, operating instructions, or recommendations. They must not replace qualified HAZOP facilitation, engineering review, or site-specific safety assessment.
\end{boundary}
\chapter{Conclusion}
On the fixed reconstructed T-301 split, typed graph reasoning is the strongest tested deterministic strategy for recovering row-derived silver associations. Its performance depends on the availability and distribution of training relationships. Stronger support and provenance-coherence filters can sharply reduce coverage; the available sample does not exercise intermediate paths; repeated holdouts show split sensitivity; and random evidence masking reduces recall and coverage. Synthetic tracks broaden the repository's data assets but remain separate from E1--E7 and retain their provenance and review limitations. The current evidence does not establish physical causal correctness or process-safety utility. Independent expert adjudication and reproducible dataset construction are required before stronger claims can be made.
\small
\bibliographystyle{plain}
\bibliography{references}
\end{document}
